% =====================================================================
\chapter{La derivada y las reglas de derivación}
\label{tema:4}
% =====================================================================

La derivada mide la tasa de cambio instantánea de una función: geométricamente, la pendiente de la recta tangente; físicamente, la velocidad. Este tema define la derivada como un límite, demuestra que la derivabilidad implica la continuidad, demuestra las reglas de derivación (suma, producto, cociente, cadena e inversa) y las aplica a las funciones elementales y a las funciones de activación. Termina con la derivación implícita y logarítmica y con las derivadas sucesivas. Se siguen los capítulos~9, 10 y~12 de Spivak~\cite{spivak}, el capítulo~4 de Apostol~\cite{apostol}, las §2.7--2.8 y el capítulo~3 de Stewart~\cite{stewart} y la unidad~1 del curso 18.01SC del MIT~\cite{mit1801}.

\section{Definición de derivada}

\begin{definicion}[Derivada]\label{def:derivada}
Sea $f$ definida en un intervalo abierto que contiene a $a$. Se dice que $f$ es \emph{derivable} en $a$ si existe y es finito el límite
\[
f'(a)=\lim_{h\to0}\frac{f(a+h)-f(a)}{h}=\lim_{x\to a}\frac{f(x)-f(a)}{x-a}.
\]
El número $f'(a)$ es la \emph{derivada} de $f$ en $a$, y la recta $y=f(a)+f'(a)(x-a)$ es la \emph{recta tangente} a la gráfica en $(a,f(a))$. La función $f'\colon x\mapsto f'(x)$, definida donde $f$ es derivable, es la \emph{función derivada}; también se escribe $\dfrac{\dd f}{\dd x}$.
\end{definicion}

El cociente $\frac{f(a+h)-f(a)}{h}$ es la pendiente de la recta secante que une $(a,f(a))$ con $(a+h,f(a+h))$; la derivada es el límite de esas pendientes cuando el segundo punto se acerca al primero. Las \emph{derivadas laterales} $f'_{+}(a)$ y $f'_{-}(a)$ se definen con los límites laterales $h\to0^{+}$ y $h\to0^{-}$; $f$ es derivable en $a$ si y solo si ambas existen y coinciden.

\addtocounter{definicion}{1}%

\begin{teorema}[Derivable implica continua]\label{teo:derivable-continua}
Si $f$ es derivable en $a$, entonces $f$ es continua en $a$.
\end{teorema}

El recíproco es falso: $\abs{x}$ y $\ReLU$ son continuas en~0 y no derivables. La siguiente reformulación de la derivada es la que mejor se generaliza a varias variables (módulo~4) y simplifica las demostraciones.

\begin{proposicion}[Derivada como aproximación lineal]\label{prop:caratheodory}
Son equivalentes: (i) $f$ es derivable en $a$ con derivada $L$; (ii) existe una función $\varphi$, continua en $a$ con $\varphi(a)=L$, tal que $f(x)=f(a)+\varphi(x)(x-a)$ para todo $x$ del dominio; (iii) $f(a+h)=f(a)+Lh+r(h)$ con $r(h)/h\to0$ cuando $h\to0$.
\end{proposicion}

La forma (iii) dice que, cerca de $a$, $f$ se parece a la función afín $f(a)+f'(a)(x-a)$ con un error que es pequeño \emph{incluso comparado con} $h$; esta \emph{aproximación lineal} es el primer escalón del polinomio de Taylor (\cref{tema:6}).

\section{Reglas de derivación}

\begin{teorema}[Suma, producto y cociente]\label{teo:reglas}
Si $f$ y $g$ son derivables en $a$, entonces:
\begin{enumerate}
  \item $(f+g)'(a)=f'(a)+g'(a)$ y $(cf)'(a)=cf'(a)$;
  \item (regla del producto) $(fg)'(a)=f'(a)g(a)+f(a)g'(a)$;
  \item (regla del cociente) si $g(a)\ne0$, $\Bigl(\dfrac fg\Bigr)'(a)=\dfrac{f'(a)g(a)-f(a)g'(a)}{g(a)^{2}}$.
\end{enumerate}
\end{teorema}

\begin{proposicion}[Potencias enteras]\label{prop:potencias}
Para todo $n\in\Z$, $(x^{n})'=nx^{n-1}$ (con $x\ne0$ si $n<0$).
\end{proposicion}

\begin{teorema}[Regla de la cadena]\label{teo:cadena}
Si $g$ es derivable en $a$ y $f$ es derivable en $g(a)$, entonces $f\circ g$ es derivable en $a$ y
\[
(f\circ g)'(a)=f'\bigl(g(a)\bigr)\,g'(a).
\]
\end{teorema}

La demostración «ingenua», que multiplica y divide por $g(x)-g(a)$, falla si $g(x)=g(a)$ para valores de $x$ arbitrariamente próximos a $a$ (por ejemplo, $g(x)=x^{2}\sin(1/x)$ en $a=0$); la función $\psi$ evita esa división. En notación de Leibniz, si $y=f(u)$ y $u=g(x)$, la regla se escribe $\dfrac{\dd y}{\dd x}=\dfrac{\dd y}{\dd u}\dfrac{\dd u}{\dd x}$.

\begin{teorema}[Derivada de la función inversa]\label{teo:inversa}
Sea $f$ continua y estrictamente monótona en un intervalo abierto $I$, derivable en $a\in I$ con $f'(a)\ne0$. Entonces $f^{-1}$ es derivable en $b=f(a)$ y
\[
(f^{-1})'(b)=\frac{1}{f'(a)}=\frac{1}{f'\bigl(f^{-1}(b)\bigr)}.
\]
\end{teorema}

Geométricamente, la gráfica de $f^{-1}$ es la simétrica de la de $f$ respecto de $y=x$, y la simetría transforma una recta de pendiente $m$ en otra de pendiente $1/m$. Si $f'(a)=0$, la tangente a $f^{-1}$ es vertical: $f(x)=x^{3}$ tiene $f'(0)=0$ y $\sqrt[3]{y}$ no es derivable en~0.

\section{Derivadas de las funciones elementales}

\begin{proposicion}[Tabla de derivadas]\label{prop:tabla}
En los dominios indicados:
\begin{enumerate}
  \item $(e^{x})'=e^{x}$ y $(\ln x)'=1/x$ para $x>0$; $(a^{x})'=a^{x}\ln a$ para $a>0$;
  \item $(x^{r})'=rx^{r-1}$ para $x>0$ y $r\in\R$;
  \item $(\sin x)'=\cos x$, $(\cos x)'=-\sin x$, $(\tan x)'=1+\tan^{2}x=1/\cos^{2}x$;
  \item $(\arcsin x)'=1/\sqrt{1-x^{2}}$ en $(-1,1)$ y $(\arctan x)'=1/(1+x^{2})$;
  \item $(\sinh x)'=\cosh x$, $(\cosh x)'=\sinh x$ y $(\tanh x)'=1-\tanh^{2}x$.
\end{enumerate}
\end{proposicion}

\begin{proposicion}[Derivadas de las funciones de activación]\label{prop:derivadas-activaciones}
Para todo $x\in\R$:
\[
\sigma'(x)=\sigma(x)\bigl(1-\sigma(x)\bigr),\qquad \softplus'(x)=\sigma(x),\qquad \tanh'(x)=1-\tanh^{2}x.
\]
En particular, $\sigma'(0)=\tfrac14$, $\tanh'(0)=1$, $0<\sigma'(x)\le\tfrac14$ y $0<\tanh'(x)\le1$.
\end{proposicion}

Que la derivada se exprese mediante el propio valor de la función es muy conveniente al programar la retropropagación: el valor $\sigma(x)$ calculado en la pasada hacia delante basta para obtener $\sigma'(x)$ en la pasada hacia atrás.

\section{Derivación implícita y logarítmica}

Una ecuación $F(x,y)=0$ puede definir $y$ como función de $x$ cerca de un punto sin que sea posible (o cómodo) despejarla. Si se supone que $y=y(x)$ es derivable, basta derivar la ecuación respecto de $x$ con la regla de la cadena y despejar $y'$. La justificación de que tal función $y(x)$ existe y es derivable es el teorema de la función implícita, que se estudia en varias variables.

\addtocounter{definicion}{1}%

\addtocounter{definicion}{1}%

\section{Derivadas sucesivas}

Si $f'$ es a su vez derivable, su derivada es la \emph{derivada segunda} $f''=(f')'$; en general, $f^{(n)}=(f^{(n-1)})'$. Se dice que $f$ es \emph{de clase $C^{n}$} en un intervalo si $f^{(n)}$ existe y es continua en él, y de clase $C^{\infty}$ si lo es para todo $n$. Los polinomios, $e^{x}$, $\sin$, $\cos$, $\sigma$ y $\tanh$ son de clase $C^{\infty}$; la $\ReLU$ es continua pero no de clase $C^{1}$ en~$\R$.

\addtocounter{definicion}{1}%

\addtocounter{definicion}{1}%

\addtocounter{definicion}{1}%

